

\def\stat{zac-ion}



\def\tit{О ПРИМЕНЕНИИ ЭКСПЕРТНЫХ МЕТОДОВ ПРИ~ОЦЕНКЕ~ЭФФЕКТИВНОСТИ 

И~КАЧЕСТВА ИНФОРМАЦИОННЫХ~СИСТЕМ}



\def\titkol{О применении экспертных методов при~оценке эффективности 

и~качества ИС} %информационных систем}



\def\aut{А.\,А.~Зацаринный$^1$, Ю.\,С.~Ионенков$^2$}



\def\autkol{А.\,А.~Зацаринный, Ю.\,С.~Ионенков}



\titel{\tit}{\aut}{\autkol}{\titkol}



\index{Зацаринный А.\,А.}

\index{Ионенков Ю.\,С.}

\index{Zatsarinny A.\,A.}

\index{Ionenkov Yu.\,S.}



%{\renewcommand{\thefootnote}{\fnsymbol{footnote}}

%\footnotetext[1]

%{Статья публикуется по представлению программного комитета VI Международной научно-

%образовательных и~ведомственных информационных сетей>>.}} 



\renewcommand{\thefootnote}{\arabic{footnote}}

\footnotetext[1]{Федеральный исследовательский центр <<Информатика и~управление>> Российской академии наук, 

AZatsarinny@ipiran.ru}

\footnotetext[2]{Федеральный исследовательский центр <<Информатика и~управление>> Российской академии наук, 

UIonenkov@ipiran.ru}



 \label{st\stat}



 \thispagestyle{headings}

 

 \vspace*{-12pt} 



 

 



  \Abst{Статья посвящена вопросу применения экспертных методов при оценке 

эффективности и~качества информационных систем (ИС). Рассмотрены общие вопросы 

применения экспертных оценок, включая проведение процедуры экспертного опроса 

и~установление согласованности мнений экспертов. Представлены различные подходы 

к~определению числа экспертов, включая формулы для проведения расчетов, и~сделаны 

сравнительные расчеты по этим формулам. Изложен формализованный подход к~оценке 

компетентности экспертов, включающий их научную квалификацию, аргументированность 

оценок и~степень знакомства с~оцениваемой проблемой.}

  

  \KW{экспертные методы; эффективность; качество; число экспертов; квалификация 

экспертов; информационная система }

  

\DOI{10.14357\!/\!08696527220205} 



\vskip 10pt plus 9pt minus 6pt



 \thispagestyle{headings}

 

 %\vspace*{-20pt}

  

\section{Введение}



  В настоящее время ИС во многом определяют 

на\-уч\-но-тех\-ни\-че\-ский потенциал страны, уровень развития ее народного 

хозяйства, образ жизни и~деятельности человека. При этом ИС относятся 

к~классу сложных больших систем с~длительным сроком эксплуатации. Одним 

из важных направлений их развития и~совершенствования является разработка 

методического аппарата, позволяющего проводить оценку их технического 

уровня на всех стадиях жизненного цикла: от замысла до снятия с~производства 

и~утилизации.

  

  Технический уровень ИС может характеризоваться двумя понятиями~--- 

<<эффективность>> и~<<качество>>. Отметим соотношение понятий 

<<качество>> и~<<эффективность>>: если эффективность характеризует 

степень соответствия системы назначению и~приспособленности к~достижению 

целей, поставленных при ее создании, то качество представляется как 

совокупность свойств системы, обусловливающих ее пригодность для 

использования по назначению~[1, 2].

  

  Анализ и~условия применения существующего методического аппарата 

оценки эффективности и~качества ИС, а также предложения по его 

совершенствованию представлены в~работах~[3--7].

  %

  В методиках, представленных в~этих работах, широко используются методы 

экспертных оценок. Данные методы основываются на знаниях специалистов 

и~накопленном ими опыте при проведении на\-уч\-но-ис\-сле\-до\-ва\-тель\-ских 

и~опыт\-но-кон\-ст\-рук\-тор\-ских работ, изучении на\-уч\-но-тех\-ни\-че\-ской 

литературы, опыте на\-уч\-но-ис\-сле\-до\-ва\-тель\-ской и~производственной 

деятельности.

  

  В литературе достаточно подробно представлены процедуры методов 

экспертных оценок, включая вопросы формирования экспертных групп, форм 

работы с~экспертами, численности и~компетентности экспертов, обработки 

экспертной информации. Основные факторы, определяющие точность 

экспертных оценок,~--- это число экспертов в~экспертной группе и~их 

компетентность.

  

  В настоящей статье рассматриваются общие вопросы применения методов 

экспертных оценок, проведено сравнение различных подходов к~определению 

требуемого числа экспертов, а также предложен методический подход к~оценке 

их компетентности.

  

\section{Методические подходы к~экспертным оценкам}

  

  Методы экспертных оценок используются, как правило, при исследовании 

объектов, анализ развития которых не поддается формализации, т.\,е.\ для 

которых трудно разработать аналитическую модель. При оценке 

эффективности и~качества ИС эти методы применяются для определения 

весовых коэффициентов и~отдельных показателей эффективности.

  

  Экспертная оценка включает~\cite{8-zac}:

  \begin{itemize}

\item составление таблицы экспертного опроса;

  \item подбор экспертов;

  \item проведение экспертизы; 

  \item анализ результатов.

  \end{itemize}

  

  Качественный и~количественный состав экспертной комиссии должен 

формироваться с~учетом широты проблемы, достоверности оценок, затрат 

ресурсов и~компетентности экспертов. 

  

  Формирование коллективного мнения экспертов осуществляется различными 

методами. Наиболее распространен из них групповой метод, когда все эксперты 

работают одновременно в~составе ка\-ко\-го-ли\-бо совета, комиссии и~т.\,п. 

Недостатки данного метода: неодинаковая активность экспертов, влияние 

авторитетов, приспособление к~мнению большинства.

  

  Наряду с~групповым методом применяется индивидуальный метод, когда 

эксперты работают отдельно и~независимо. Обобщенное мнение формируется 

путем статистической обработки мнений отдельных экспертов. В~этом случае 

удается исключить недостатки, присущие групповому методу, но появляется 

опасность нивелирования мнений отдельных экспертов, следования за 

обезличенным статистическим большинством.

  

  Успех экспертного опроса зависит от числа исходных показателей: чем 

больше показателей подлежит оценке, тем большей совокупности показателей 

эксперты присваивают одинаковую важность. Иначе говоря, с~возрастанием 

числа показателей эксперты становятся все более безразличными в~своих 

суждениях. Поскольку в~общем случае число исходных показателей не может 

быть сокращено без ущерба для полноты описания системы, полезно введение 

обобщенных показателей, объединяющих ряд исходных показателей, и~их 

группировка. 

  

  Процедура экспертного опроса подразделяется на три стадии.

  

  На первой стадии составляется таблица экспертного опроса и~приложение 

к~ней. В~таб\-ли\-цу заносятся характеристики и~показатели систем, по 

которым необходимо выяснить мнение специалистов. Каждому специалисту 

предлагается оценить их важность в~баллах, а~также указать свою 

квалификацию, степень знакомства с~рассматриваемыми системами, источники 

аргументации.

  

  На второй стадии определяется круг экспертов и~проводится 

непосредственно экспертный опрос.

  

  На третьей стадии проводится обработка данных: вычисляются средние 

оценки и~устанавливается степень согласованности мнений экспертов. 

  

  Сначала оценки, данные каждым экспертом, нормируются по формуле

  \begin{equation*}

  \beta_{kn} = \fr{C_{kn}}{\sum\nolimits_{n=1}^N C_{kn}}\,,

 % \label{e1-zac}

  \end{equation*}

где $C_{kn}$~--- оценка $k$-м экспертом важности $n$-го показателя 

(в~баллах); $N$~--- число показателей.

  

  Затем вычисляются средние значения оценок каждого показателя

  \begin{equation*}

  \overline{\beta}_b=\fr{1}{K} \sum\limits^K_{k=1} \beta_{kn}

 % \label{e2-zac}

  \end{equation*}

и соответствующие среднеквадратические отклонения 

\begin{equation*}

\sigma= \sqrt{\fr{1}{K} \sum\limits^K_{k=1} \left( \beta_{kn}-

\overline{\beta}_n\right)^2}\,,

%\label{e3-zac}

\end{equation*}

где $K$~--- число экспертов. 



После этого определяются коэффициенты вариации

\begin{equation*}

\upsilon_n=\fr{\sigma_n}{\overline{\beta}_n}\,,

%\label{e4-zac}

\end{equation*}

характеризующие степень согласованности мнений экспертов о важности 

отдельных показателей. Чем больше величина~$\upsilon_n$, тем меньше 

согласованность экспертов в~оценке $n$-го показателя; при одинаковом 

(полностью согласованном) мнении экспертов эта величина равна нулю. 



  Степень согласованности мнений экспертов по всем показателям 

оценивается коэффициентом конкордации Кенделла~\cite{8-zac}:

  \begin{equation*}

  \omega= \fr{12\sum\nolimits^N_{n=1} \left( S_n-\overline{S}\right)^2} 

{K^2\left( N^3-N\right) -K\sum\nolimits^K_{k=1} T_k}\,.

  %\label{e5-zac}

  \end{equation*}

Здесь  $S_n=\sum\nolimits^K_{k=1} R_{kn}$~--- сумма рангов, присвоенных 

экспертами $n$-му показателю, где $R_{kn}$~--- ранг $n$-го показателя, 

назначенный $k$-м экспертом; $\overline{S}\hm= (1/N)\sum\nolimits^N_{n=1} 

S_n$~--- среднее значение сумм рангов;

$$

T_k= \sum\limits_{m=1}^{M_k} \left( t^3_{km}-t_{km}\right)\,,

$$

где $t_{km}$~--- число одинаковых рангов $m$-го типа в~оценках $k$-го 

эксперта, $M_k$~--- количество групп показателей с~совпавшими рангами 

в~оценках $k$-го эксперта.



  Коэффициент конкордации изменяется в~пределах от нуля до единицы, 

причем $\omega\hm=1$ соответствует полной согласованности мнений 

экспертов.

  

  При обработке данных по группам экспертов вычисляется коэффициент 

ранговой корреляции Спирмэна~\cite{8-zac}

  \begin{equation*}

  r=1-\fr{6\sum\nolimits^N_{n-1} \left( R_{ni}-R_{nj}\right)^2} {N\left( N^2-

1\right)}\,,

 % \label{e6-zac}

  \end{equation*}

показывающий связь между оценками, назначенными экспертами $i$-й и~$j$-й 

групп. 



  Если согласованность мнений экспертов высокая и~неслучайная 

(коэффициенты конкордации и~ранговой корреляции значимы и~близки 

к~единице), то найденные весовые коэффициенты и~отдельные показатели 

могут быть использованы для расчета обобщенных оценок эффективности.



\section{Методические подходы к~определению числа экспертов}



  При проведении экспертных оценок очень важным вопросом является 

определение числа экспертов. Число экспертов должно быть достаточно 

большим, чтобы отдельные мнения не имели неправомерно большое значение, 

однако при очень большом числе экспертов снижается уровень их 

компетентности, что может привести к~снижению точности экспертных оценок.

  

  К экспертам предъявляются следующие требования: глубокие знания 

в~оцениваемой области; наличие научного интереса к~оцениваемым аспектам 

проблемы; наличие производственного или исследовательского опыта по 

оцениваемой проблеме.

  

  В научно-технической литературе предлагается ряд подходов к~определению 

требуемого числа экспертов, базирующихся на статистических методах.

  

  В частности, в~\cite{9-zac} предложена следующая формула для определения 

тре\-бу\-емо\-го числа экспертов:

  \begin{equation}

  n=\fr{t^2\sigma^2 N}{\Delta^2N+t^2\sigma^2}\,,

  \label{e7-zac}

  \end{equation}

где $n$~--- число экспертов;

$t$~--- нормированное отклонение, соответствующее принятому уровню 

доверительной вероятности;

$\Delta$~--- предельная ошибка выборки;

$\sigma^2$~--- дисперсия исследуемого признака в~генеральной совокупности;

$N$~--- объем генеральной совокупности.

  

  Если дисперсия исследуемого признака даже приблизительно не известна, то 

она принимается равной своему максимуму~--- 0,25 (0,5$\times$0,5) 

и~формула~(\ref{e7-zac}) приобретает вид:

  \begin{equation}

  n=\fr{0{,}25 t^2N}{\Delta^2 N+0{,}25 t^2}\,.

  \label{e8-zac}

  \end{equation}

  

  В~\cite{10-zac} для определения числа экспертов использовано соотношение, 

которое применяется для вычисления погрешности наблюдений:

  \begin{equation}

  n=\fr{t_p^2}{\varepsilon_1^2}\,.

  \label{e9-zac}

  \end{equation}

Здесь $n$~--- число экспертов;

$\varepsilon_1\hm= \varepsilon/S$~--- предельно допустимая 

относительная ошибка экспертной оценки, где

$S$~--- среднеквадратичное отклонение распределения оценок;

$t_p$~--- коэффициент Стьюдента, определяющий ширину доверительного 

интервала и~зависимость от величины вероятности оценки~${\sf P}$ ($t_p$~--- 

табулированная величина). 

  

  В зависимости от заданной погрешности экспертной оценки и~выбранной 

величины вероятности~$P$ может быть определена возможная численность 

экспертов (табл.~1)~\cite{10-zac}.

  

\begin{table}[b]\small %tabl1

\vspace*{-9pt}

\begin{center}

\Caption{Минимально допустимое число экспертов в~группе}

\vspace*{2ex}



\tabcolsep=7pt

\begin{tabular}{|c|c|c|c|c|c|c|c|c|}

\hline

  \multicolumn{1}{|c|}{\raisebox{-6pt}[0pt][0pt]{$\varepsilon_1$}}&\multicolumn{8}{c|}{Вероятность оценки~${\sf P}$}\\

  \cline{2-9}

  &0,99&0,95&0,90&0,85&0,8&0,75&0,7&0,65\\

  \hline

  0,5&26&15&11&\hphantom{9}8&\hphantom{9}7&\hphantom{9}5&\hphantom{9}4&\hphantom{9}4\\

  0,3&74&43&31&23&19&15&12&10\\

  \hline

  \end{tabular}

  \end{center}

  %\vspace*{-9pt}

  \end{table}

  

  В~\cite{11-zac} минимальное число экспертов предлагается определять по 

формуле:

  \begin{equation}

  n=0{,}5 \left( \fr{3}{\alpha}+5\right)\,,

  \label{e10-zac}

  \end{equation}

где $0<\alpha \leq 1$~--- параметр, задающий минимальный уровень ошибки 

экспертизы.

  

  В соответствии с~\cite{12-zac}, для проведения экспертной оценки 

необходимо привлечение не менее 7--9 экспертов.

  

  Для сравнения этих подходов и~определения требуемого числа экспертов 

проведем расчеты по представленным выше  

формулам~(\ref{e8-zac})--(\ref{e10-zac}).

  

  При расчете по формуле~(\ref{e8-zac}) примем величину доверительной 

вероятности, представляющей собой табулированную величину, равной~0,9545 

($t\hm=2$), при этом $\Delta= t \sqrt{\sigma^2/N}\hm= 2 \cdot 0{,}1235 \hm= 

0{,}247$, а~объем генеральной совокупности, анализируемой экспертами, $N$ 

примем равным~1500. Тогда число экспертов составит

  $$

  n= \fr{0{,}25 \cdot 1500\cdot 2^2}{0{,}247^2\cdot 1500 +0{,}25\cdot 

2^2}=16\,.

  $$

  

  При использовании формулы~(\ref{e9-zac}) для предельно допустимой 

относительной ошибки экспертной оценки, равной~0,3, число экспертов при 

вероятности оценки~0,8 составит 19~чел., а при вероятности оценки  

0,75~--- 15~чел. 

  

  При расчете по формуле~(\ref{e10-zac}) при $\alpha\hm = 0{,}2$ получаем 

10~чел., а~при $\alpha\hm = 0{,}3$~---  8~чел.

  

  Таким образом, в~соответствии с~проведенными по разным формулам 

расчетами, для достоверного результата экспертного оценивания (вероятность 

оценки не менее~0,8) число экспертов должно составлять 10--15~чел., а~для 

получения более высокой достоверности оценок число экспертов должно быть 

увеличено.

  

  Кроме того, достоверность оценок, полученных в~результате экспертного 

опроса, существенно зависит от компетентности привлекаемых экспертов. 

  

\section{Методический подход к~оценке компетентности экспертов}



  Не менее важным вопросом, чем определение числа экспертов в~группе, 

является оценка компетентности экспертов. Следует отметить, что объективные 

оценки экспертов на практике трудно реализуемы и~в основном используются 

субъективные способы: оценка на основании документов о квалификации, 

взаимооценка и~самооценка.

  

  Для оценки компетентности экспертов предлагается доработанный авторами 

статьи формализованный подход~\cite{10-zac, 13-zac}, учитывающий научную 

квалификацию эксперта ($\mathrm{К}_{\mathrm{н}}$), аргументированность его 

оценок ($\mathrm{К_а}$) и~степень знакомства с~оцениваемой проблемой 

($\mathrm{К_з}$).

  

  Научная квалификация эксперта ($\mathrm{К_н}$) оценивается в~соответствии 

с~табл.~2~\cite{10-zac, 13-zac}. 

  

  

\begin{table}\small %tabl2

\begin{center}

\Caption{Примерный уровень научной квалификации эксперта}

\vspace*{2ex}



\tabcolsep=4pt

\begin{tabular}{|l|c|c|c|c|}

\hline

&\multicolumn{4}{c|}{Значения коэффициента 

квалификации $\mathrm{К_н}$}\\

\cline{2-5}

\multicolumn{1}{|c|}{\tabcolsep=0pt\begin{tabular}{c}Должность\\ \  \end{tabular}}&

\tabcolsep=0pt\begin{tabular}{c}Без ученой\\ степени\end{tabular}&

\tabcolsep=0pt\begin{tabular}{c}Кандидат\\ наук\end{tabular}&

\tabcolsep=0pt\begin{tabular}{c}Доктор\\ наук\end{tabular}&

\tabcolsep=0pt\begin{tabular}{c}Член-корр.,\\ академик\end{tabular}\\

\hline

Старший научный сотрудник&0,15&\hphantom{99}0,225&0,3&0,5\\

\hline

\tabcolsep=0pt\begin{tabular}{l}Начальник лаборатории,\\ руководитель группы\end{tabular}&

0,2\hphantom{9}&0,3&0,4&0,6\\

\hline

\tabcolsep=0pt\begin{tabular}{l}Начальник отдела,\\ заместитель начальника отдела\end{tabular}&0,25&

\hphantom{99}0,375&0,5&\hphantom{9}0,75\\

\hline

\tabcolsep=0pt\begin{tabular}{l}Руководитель комплекса,\\ заместитель руководителя комплекса\end{tabular}&

0,3\hphantom{9}&0,5&0,6&0,9\\

\hline

\tabcolsep=0pt\begin{tabular}{l}Директор, заместитель директора,\\ научный руководитель темы\end{tabular}&

0,4\hphantom{9}&0,6&0,8&1,0\\

\hline

\end{tabular}

\end{center}

  \vspace*{-9pt}

\end{table}



  

  В работах~\cite{8-zac, 10-zac} предложен ряд показателей для 

аргументированности оценок эксперта. Но они не в~полной мере учитывают 

состояние и~перспективы развития современных ИС. 

  

  

  Исходя из этого, авторами предлагается оценивать аргументированность 

оценок эксперта ($\mathrm{К}_{\mathrm{а}}$) по следующим показателям: стаж 

работы в~данной области; производственный  

и~на\-уч\-но-ис\-сле\-до\-ва\-тель\-ский опыт; проведенный теоретический 

анализ проблемы; учет отечественных и~зарубежных публикаций в~данной 

области; учет современных тенденций развития ИС  

и~ин\-фор\-ма\-ци\-он\-но-те\-ле\-ком\-му\-ни\-ка\-ци\-он\-ных технологий; 

наличие научных трудов; личностные характеристики эксперта 

(эмоциональность и~т.\,д.). Каждая из этих составляющих оценивается 

соответствующим коэффициентом в~диапазоне от~0 до~1, при этом суммарный 

коэффициент аргументации не должен превышать единицы. Оценка 

аргументации мнения эксперта представлена в~табл.~3.



\begin{table}\small %tabl3

\begin{center}

\Caption{Оценка аргументации мнения эксперта}

\vspace*{2ex}



\begin{tabular}{|l|c|}

\hline

\multicolumn{1}{|c|}{Источники аргументации}&\tabcolsep=0pt\begin{tabular}{c}Коэффициент\\ аргументации\end{tabular}\\

\hline

Стаж работы &0,15\\

Производственный и~научно-исследовательский опыт&0,2\hphantom{9}\\

Теоретический анализ проблемы&0,15\\

Учет отечественных и~зарубежных публикаций&0,15\\

Учет современных тенденций развития ИС&0,2\hphantom{9}\\

Наличие научных трудов&0,1\hphantom{9}\\

Личностные характеристики эксперта&0,05\\

\hline

\end{tabular}

\end{center}

  \vspace*{-9pt}

\end{table}

  

  Степень знакомства эксперта с~оцениваемой проблемой 

($\mathrm{К}_{\mathrm{з}}$) характеризуется выставляемым экспертом баллом 

самооценки (от~0 до~10), причем максимальному баллу~(10) соответствует 

полное знакомство с~оцениваемой проблемой, а~минимальному баллу~(0)~--- 

полное отсутствие знакомства. Затем проводится нормирование оценки с~ее 

переводом в~диапазон от~0 до~1.

  

  Комплексный показатель компетентности $i$-го эксперта в~простейшем 

случае, принимая одинаковые значения отдельных составляющих, можно 

вычислить по формуле:



\vspace*{2pt}



\noindent

  \begin{equation*}

  \mathrm{К}_i = \fr{\mathrm{К_н + К_а + К_з}}{3}\,.

  % \label{e11-zac}

   \end{equation*}

   

   \vspace*{-3pt}

  

  В некоторых работах (см., например,~\cite{14-zac, 15-zac}) предлагается более 

сложный подход, связанный с~вычислением весовых коэффициентов 

отдельных элементов комплексного показателя эффективности 

с~использованием метода анализа иерархий. 

  

  Комплексный показатель компетентности эксперта характеризует весомость 

мнения данного эксперта и~показывает, с~каким весом следует включать 

количественные оценки данного эксперта по рассматриваемой проблеме 

в~статистическую обработку.



\vspace*{-4pt}



\section{Заключение}



\vspace*{-1pt}



  

  При оценке эффективности и~качества ИС широко используются методы 

экспертных оценок. Эти методы применяются для решения задач, которые не 

могут быть решены формализованными методами из-за отсутствия 

необходимых данных по характеристикам создаваемой системы. В~этом случае 

возникает необходимость использования экспертных оценок, основанных на 

опыте и~знаниях специалистов в~исследуемой об\-ласти.

  

  Процедура проведения экспертных оценок получила широкое развитие, она 

всесторонне обоснована и~максимально формализована. Кроме того, хорошо 

разработан математический аппарат оценки точности экспертных оценок. 

Вместе с~тем одной из основных проблем при проведении экспертиз остается 

определение числа экспертов и~их компетентности. 

  

  

  В данной статье рассмотрен ряд подходов к~определению числа экспертов, 

проведены расчеты по различным формулам и~сделан вывод о том, что число 

экспертов при вероятности оценки не менее~0,8 должно составлять  

10--15~человек, а для получения более высокой достоверности оценок число 

экспертов должно быть увеличено.

  

  Для оценки компетентности экспертов предложен формализованный подход, 

учитывающий научную квалификацию эксперта, аргументированность его 

оценок и~степень знакомства с~оцениваемой проблемой. Даны предложения по 

расчету комплексного показателя компетентности экспертов.



\vspace*{-4pt}



{\small\frenchspacing

{%\baselineskip=10.8pt

\addcontentsline{toc}{section}{Литература}

\begin{thebibliography}{99}



\vspace*{-1pt}



\bibitem{2-zac}

ГОСТ 15467-79. Управление качеством продукции. Основные понятия. Термины 

и~определения.~--- М.: Стандартинформ, 2009. 21~с.



\bibitem{1-zac}

ГОСТ 34.003-90. Информационная технология. Автоматизированные системы. Термины 

и~определения.~--- М.: Стандартинформ, 2005. 14~с.



 \bibitem{3-zac}

\Au{Зацаринный А.\,А., Ионенков~Ю.\,С., Шабанов~А.\,П.} К~вопросу о сравнительной 

оценке эффективности ситуационных центров~/\!\!/ Системы и~средства информатики, 2013. 

Т.~23. №\,2. С.~155--171.

\bibitem{4-zac}

\Au{Зацаринный А.\,А., Ионенков~Ю.\,С.} К~вопросу оценки эффективности 

автоматизированных систем с~использованием метода анализа иерархий~/\!\!/ Системы 

и~средства информатики, 2015. Т.~25. №\,3. С.~161--178.

\bibitem{5-zac}

\Au{Зацаринный А.\,А., Ионенков~Ю.\,С.} Некоторые аспекты оценки эффективности 

автоматизированных информационных систем на различных стадиях их жизненного 

цикла~/\!\!/ Системы и~средства информатики, 2016. Т.~26. №\,3. С.~123--136.

\bibitem{6-zac}

\Au{Зацаринный А.\,А., Ионенков Ю.\,С., Сучков~А.\,П.} Некоторые аспекты оценки 

эффективности облачных технологий~/\!\!/ Системы и~средства информатики, 2018. Т.~28. 

№\,3. С.~104--117.

\bibitem{7-zac}

\Au{Зацаринный А.\,А., Ионенков Ю.\,С.} Некоторые методические аспекты выбора 

показателей эффективности информационных систем~/\!\!/ Системы высокой доступности, 

2019. №\,4. С.~19--26.

\bibitem{8-zac}

\Au{Окунев Ю.\,Б., Плотников В.\,Г.} Принципы системного подхода к~проектированию 

в~технике связи.~--- М.: Связь, 1976. 183~с.

\bibitem{9-zac}

\Au{Кузовкова Т.\,А., Салютина Т.\,Ю., Шаравова~О.\,И.} Статистика инфокоммуникаций.~--- 

М.: Горячая ли\-ния-Те\-ле\-ком, 2019. 548~с. 

\bibitem{10-zac}

\Au{Семенов С.\,С.} Оценка качества и~технического уровня сложных систем: Практика 

применения метода экспертных оценок.~--- М.: ЛЕНАНД, 2015. 352~с.

\bibitem{11-zac}

\Au{Литвак Б.\,Г.} Экспертные оценки и~принятие решений.~--- М.: Патент, 1996. 298~с.

\bibitem{12-zac}

\Au{Саати Т.} Принятие решений: Метод анализа иерархий.~--- М.: Радио и~связь, 1993. 

278~с.

\bibitem{13-zac}

\Au{Бобровников Г.\,Н., Клебанов А.\,И.} Комплексное прогнозирование создания новой 

техники.~--- М.: Экономика, 1989. 205~с.

\bibitem{14-zac}

\Au{Чернышева Т.\,Ю.} Иерархическая модель оценки и~отбора экспертов~/\!\!/ Доклады 

ТУСУР, 2009. №\,1(19). Ч.~1. С.~168--173.

\bibitem{15-zac}

\Au{Марычева П.\,Г.} Методика оценки компетентности экспертов~/\!\!/ Вестник Самарского 

государственного технического университета. Сер.: Технические науки, 2018. №\,4(60). 

С.~29--40. 

       \end{thebibliography}

} }







\vspace*{-6pt}



\hfill{\small\textit{Поступила в~редакцию 28.12.21}}



\vspace*{8pt}



\hrule



\vspace*{2pt}



%\newpage



%\vspace*{-28pt}



\hrule



%\vspace*{8pt}



\def\tit{ON THE USE OF EXPERT METHODS IN~EVALUATING EFFECTIVENESS AND~QUALITY 

OF~INFORMATION SYSTEMS}







\def\titkol{On the use of expert methods in~evaluating effectiveness and~quality 

of~information systems}



\def\aut{A.\,A.~Zatsarinny and Yu.\,S.~Ionenkov}



\def\autkol{A.\,A.~Zatsarinny and Yu.\,S.~Ionenkov}



\titel{\tit}{\aut}{\autkol}{\titkol}



\vspace*{-15pt}



 \noindent

Federal Research Center ``Computer Science and Control'' of the Russian Academy of Sciences, 

44-2~Vavilov Str., Moscow 119333, Russian Federation





\def\leftfootline{\small{\textbf{\thepage}

\hfill Sistemy i Sredstva Informatiki~---

Systems and Means of Informatics 2022 vol~32 no\,2}

}%

 \def\rightfootline{\small{Sistemy i Sredstva Informatiki~---

 Systems and Means of Informatics 2022 vol~32 no\,2

\hfill \textbf{\thepage}}}



\vspace*{3pt}





\Abste{The article is devoted to the application of expert methods in assessing effectiveness and 

quality of information systems. The general issues of the\linebreak\vspace*{-12pt}}



\Abstend{application of expert assessments, 

including the conduct of the expert survey procedure and the establishment of consistency of expert 

opinions, are considered. Various approaches to determining the number of experts, including 

formulas for calculations, are presented and comparative calculations based on these formulas are 

carried out. A formalized approach to assessing the competence of experts is presented, including 

their scientific qualifications, reasonableness of assessments, and the degree of familiarity with the 

problem being evaluated.}



\KWE{expert methods; efficiency; quality; number of experts; qualification of experts; information 

system}



\DOI{10.14357\!/\!08696527220205} 



%\vspace*{-16pt}





 %\Ack

%\noindent



 





%\vspace*{-6pt}



\renewcommand{\bibname}{\protect\rmfamily References}

%\renewcommand{\bibname}{\large\protect\rm References}



{\small\frenchspacing

{%\baselineskip=10.8pt

\addcontentsline{toc}{section}{References}

\begin{thebibliography}{99}



\bibitem{2-zac-1}

 GOST 15467-79. 2009. Upravlenie kachestvom produktsii. Osnovnye ponyatiya. Ter\-mi\-ny 

i~opredeleniya [Product quality control. Basic concepts. Terms and definitions]. Moscow: 

Standardinform Publs. 21~p.



\bibitem{1-zac-1}

GOST 34.003-90. 2005. Informatsionnaya tekhnologiya. 

 Avtomatizirovannye sistemy. Terminy i~opredeleniya [Information 

technology. Set of standards for automated systems. Automated systems. Terms and definitions]. 

 Moscow: Standardinform Publs. 14~p.

 

\bibitem{3-zac-1}

\Aue{Zatsarinnyy, A.\,A., Yu.\,S.~Ionenkov, and A.\,P.~Shabanov.} 2013. K~voprosu 

o~sravnitel'noy otsenke effektivnosti situatsionnykh tsentrov [Regarding comparative evaluation of 

situational centers efficiency]. \textit{Sistemy i~Sredstva Informatiki~--- Systems and Means of 

Informatics} 23(2):155--171.

\bibitem{4-zac-1}

\Aue{Zatsarinnyy, A.\,A., and Yu.\,S.~Ionenkov.} 2015. K~voprosu otsenki effektivnosti 

av\-to\-ma\-ti\-zi\-ro\-van\-nykh sistem s~ispol'zovaniem metoda analiza ierarkhiy [Regarding automated 

systems efficiency evaluation using analytic hierarchy process]. \textit{Sistemy i~Sredstva 

Informatiki~--- Systems and Means of Informatics} 25(3):161--178.

\bibitem{5-zac-1}

\Aue{Zatsarinnyy, A.\,A., and Y.\,S.~Ionenkov.} 2016. Nekotorye aspekty otsenki effektivnosti 

avtomatizirovannykh informatsionnykh sistem na razlichnykh stadiyakh ikh zhiznennogo tsikla [On 

aspects of automated information system efficiency evaluation at different stages of lifecycle]. 

\textit{Sistemy i~Sredstva Informatiki~--- Systems and Means of Informatics} 26(3):123--136.

\bibitem{6-zac-1}

\Aue{Zatsarinny, A.\,A., Yu.\,S.~Ionenkov, and A.\,P.~Suchkov.} 2018. Nekotorye aspekty otsen\-ki 

effektivnosti oblachnykh tekhnologiy [Some aspects of cloud computing efficiency estimation]. 

\textit{Sistemy i~Sredstva Informatiki~--- Systems and Means of Informatics} 28(3):104--117.

\bibitem{7-zac-1}

\Aue{Zatsarinny, A.\,A., and Y.\,S.~Ionenkov.} 2019. Nekotorye metodicheskie aspekty vybora 

po\-ka\-za\-te\-ley effektivnosti informatsionnykh sistem [Some methodological aspects of the choice of 

performance indicators of information systems]. \textit{Sistemy vysokoy dostupnosti} [Highly Available 

Systems] 15(4):19--26.

\bibitem{8-zac-1}

\Aue{Okunev, Y.\,B., and V.\,G.~Plotnikov.} 1976. \textit{Printsipy sistemnogo podkhoda k~proektirovaniyu 

v~tekhnike svyazi} [Systematic approach principles to design in communication technics]. Moscow: 

Svyaz'. 183~p.

\bibitem{9-zac-1}

\Aue{Kuzovkova, T.\,A., T.\,Yu.~Salyutina, and O.\,I.~Sharavova.} 2019. \textit{Statistika infokommunikatsiy} 

[Statistics of infocommunications]. Moscow: Goryachaya liniya-Telekom. 548~p.

\bibitem{10-zac-1}

\Aue{Semenov, S.\,S.} 2015. \textit{Otsenka kachestva i~tekhnicheskogo urovnya slozhnykh sistem: Praktika 

primeneniya metoda ekspertnykh otsenok} [Assessment of the quality and technical level of complex 

systems: The practice of applying the method of expert assessments]. Moscow: LENAND. 352~p.

\bibitem{11-zac-1}

\Aue{Litvak, B.\,G.} 1996. \textit{Ekspertnye otsenki i~prinyatie resheniy} [Expert assessments and 

decision-making]. Moscow: Patent. 298~p.

\bibitem{12-zac-1}

\Aue{Saati, T.} 1993. \textit{Prinyatie resheniy. Metod analiza ierarkhiy} [Decision making. Analytic 

hierarchy process]. Moscow: Radio i~svyaz'. 278~p.

\bibitem{13-zac-1}

\Aue{Bobrovnikov, G.\,N., and A.\,I.~Klebanov.} 1989. \textit{Kompleksnoe prognozirovanie so\-zda\-niya 

novoy tekhniki} [Comprehensive forecasting of the creation of new equipment]. Moscow: 

Ekonomika. 205~p.

\bibitem{14-zac-1}

\Aue{Chernysheva, T.\,Yu.} 2009. Ierarkhicheskaya model' otsenki i otbora ekspertov [Hierarchical 

model of evaluation and selection of experts]. \textit{Doklady TUSUR} [Proceedings of the TUSUR 

University] 1-1(19):168--173.

\bibitem{15-zac-1}

\Aue{Marycheva, P.\,G.} 2018. Metodika otsenki kompetentnosti ekspertov [The method of assessment 

the competence of experts]. \textit{Vestnik Samarskogo gosudarstevennogo tekh\-ni\-che\-sko\-go universiteta. Ser. Tekh\-ni\-che\-skie nauki}

[Vestnik of Samara State Technical University. Technical sciences ser.] 4(60):29--40.

\end{thebibliography}

} }



\vspace*{-3pt}



\hfill{\small\textit{Received December 28, 2021}} 



%\vspace*{-14pt} 



\Contr



\noindent

\textbf{Zatsarinny Alexander A.} (b.\ 1951)~--- Doctor of Science in technology, professor, 

principal scientist, Federal Research Center ``Computer Science and Control'' of the Russian 

Academy of Sciences, 44-2~Vavilov Str., Moscow 119333, Russian Federation; 

\mbox{AZatsarinny@ipiran.ru}



\vspace*{3pt}



\noindent

\textbf{Ionenkov Yuriy S.} (b.\ 1956)~--- Candidate of Science (PhD) in technology, senior 

scientist, Federal Research Center ``Computer Science and Control'' of the Russian Academy of 

Sciences, 44-2~Vavilov Str., Moscow 119333, Russian Federation; \mbox{uionenkov@ipiran.ru}







\label{end\stat}



\renewcommand{\bibname}{\protect\rm Литература} 

